\subsection{Polynomials in One Variable}

\begin{definition}[title = Definition 1.5.1]
    For $f = a_o + a_1 x + \ldots + a_m x^m \in k[x]\backslash \{0\}$ with $a_0, \ldots, a_m \in k$ and $a_m \neq 0$, we call $a_mx^m$ the \textit{leading term of $f$}.

    Notation: $LT(f)$
\end{definition} 

\begin{proposition}[title = Proposition 1.5.2]
  For every choice of $f \in k[x]$ and $g \in k[x]\backslash \{0\}$ there exisyts a unique decomposition \[f = q \cdot g + r\] with $q,r \in k[x]$ and 
  $\deg r < \deg g$. Furthermore, there exists an algrithm that calculates $g$ and $r$ for given $f$ and $g$.

  \paragraph{Proof:}  In our algorithm we start with 
    \[\underbrace{q =0}_{\text{ quotient in progress }} \text{ and } \underbrace{r = f}_{\text{ remainder in progress }}\] 

    We will maintain the equality $f = g \cdot q + r$ as the invariant, making the $\deg r$ smaller in each iteration until the condition $\deg r < \deg g$ is satisfied. 

    If $\deg r < \deg g$, then we can terminate. Otherwise, we want to update $r$ using $g$ to make $r$ smaller. 

    \[\deg r \geq \deg g  \implies \\
        f  = g \cdot \underbrace{q + \frac{LT(r)}{LT(g)}}_{\text{new } q} + \underbrace{r - \frac{LT(r)}{LT(g)} \cdot g}_{\text{new } r}\]
    by taking into account that $f = g \cdot q + r$. 

    This leads to: 
    \begin{algorithm}[title = Univariate Division Algorithm]
        \textbf{Assumption}: $f,g \in k[x], g \neq 0$  \\
        \textbf{Result}: quotient $q$ and the remainder $r$ of the long division of $f$ by $g$. 
\begin{lstpseudo}
r:=f    % Remainder in progress 
q:=0    % Quotient in progress 

while deg r >= deg g:
    T := LT(r)/LT(g)  
    q := q + T   % update of the quotient in progress 
    r := r - T * g  % update of the remainder in progress  
 
return (q,r)
\end{lstpseudo}
    \end{algorithm}

    The algorithm is correct and is terminating in finitely many steps, because we maintain the invariant $f = g \cdot q + r$ and the updates of $q$ and  $r$ are such 
    that $\deg r$ gets strickly smaller with each iteration. \\

    It remains to show the uniqueness of the decomposition. 

    Assume $f = q \cdot g+r = g \cdot Q + R$ with $q,r,Q,R \in k[x]$ and $\deg r < \deg g$ and $\deg R < \deg g$. 
    
    If $q = Q$, then $r = f-g \cdot q = f -g \cdot Q = R$ and so $r = R$. 

    If $q \neq Q$, then $0 = f -f = (g \cdot Q - R) - (g \cdot q -r) = g \cdot (Q -q) + r - R$ 

    \[\quad \underbrace{R - r}_{degree < \deg g} = \underbrace{g}_{\deg g} \cdot \underbrace{(Q- q)}_{\neq 0}\]
    
    The left-hand side has degree $< \deg g$, but the right-hand side has degree $\geq \deg g$, which is a contradiction.  

    \[\Longrightarrow Q = q\]
    \begin{flushright}
      $\Box $
    \end{flushright}
\end{proposition}

\begin{corollary}[title = Corollary 1.5.3]
    Every $f \in k[x]\backslash\{0\}$ has at most $\deg \ f$ roots in $k$. 

    \paragraph{Proof:} Induction on $m := \deg \ f$.  \\[0.3em]
    For $m = 0$, this is clear.  \\[0.3em]
    Let $m \geq 1$ and assume that the assertion is verified for polynomials of degree $< m$. \\[0.3em]
    Consider any $f \in k[x]\backslash\{0\}$ of $\deg \ f= m$.  
    
    If $f$ has no roots in $k$, we are done, because $0\leq m$. 

    If $f$ has at least one root $a \in k$, we can use long division to divide $f$ by $x -a$. We obtain $f = (x-a) \cdot q - r$ where $q \in k[x]$ an $\deg \ r < \deg(x-a) =1$ 
    meaning $r \in k$. 
    Since $f(a) = 0$, we have $f(a) = (a -a) \cdot q(a) + r$ 
    \begin{addmargin}[0.6cm]{0cm}  
        $\Rightarrow \quad r = 0$ \\
        $\Rightarrow \quad f = (x-a) \cdot q$, where $q$ is a polynomial of degree $m -1$. \\
        \hspace*{0.8cm}By induction, $q$ has at most $m -1$ roots in $k$. \\
        $\Rightarrow \quad f$ has at most $m$ roots (the root $a$ and the roofs of $g$) 
    \end{addmargin}

    \begin{flushright}
      $\Box $
    \end{flushright}
\end{corollary}


\begin{corollary}[title = Corollary 1.5.4]
  Every ideal $I$ in $k[x]$ has the form $I = \langle h \rangle$ for some $h \in k[x]$. 

  \paragraph{Proof:}  If $I = \{0\}$, the assertion is trivial as we can fix $h = 0$. \\[0.3em]
  Assume  $I \neq \{0\}$. Choose a polynomial $h \in I$ with $h \neq 0$ with the lowest possible degree. We claim that $I = \langle h \rangle$. The inclusion $I \supseteq \rangle h \rangle$ 
  is clear because $h \in I$. \\[0.5em] 
  Let's show the converse: $I \subseteq \langle h \rangle$.  

  Consider any $f \in I$. We can use long division to divide $f$ by $h$: $f = h \cdot q + r$ with $q,r \in k[x]$ and $\deg r < \deg h$. 
  But then \[r = \underbrace{f}_{\in I} - \underbrace{h}_{\in I} \cdot q \in I\] has a smaller degree than h, contradicting the choice of $h$ whwnever $r \neq 0$. So, $r = 0$. 
  \[
    \Rightarrow f = h \cdot q + 0 = h \cdot q \in \langle h \rangle
  \]

    \begin{flushright}
      $\Box $
    \end{flushright}
\end{corollary} 

\begin{definition}[title = Definition 1.5.5]
  A \textit{principal ideal} of a ring is an ideal that can be generated by just one element of the ring. A principal ideal domain (PID) is a commutative ring, in which every 
  ideal is principal. 
\end{definition}

We have shown $k[x]$ is a PID. Another example is $\mathbb{Z}$. 

\begin{definition}[title = Definition 1.5.6 (Greatest Common Divisor (GCD))]
    $f, g \in k[x]$ is a polynomial $h \in k[x]$ such that: 
    \begin{itemize}
      \item[(i)] $h$ divides both $f$ and $g$ 
      \item[(ii)] every $p \in k[x]$ that divides $f$ and $g$, also divides $h$. 
    \end{itemize}
    Nottion: $h = \gcd(f,g)$
\end{definition}

The notatoin suggests that $gcd(f,g)$ is unique, but we will see that $gcd(f,g)$ is unique up to constant non-zero multiple only. Thus, by writing $gcd(g,f)$ we abuse 
notation (but we won't missuse is hopefully).

\clearpage 

\begin{tcolorbox}[sharp corners, title = In Math:, colback = airforceblue!5!white, colframe = airforceblue!95!black, fonttitle = \bfseries]
   \textbf{Abuse of notation:} tentatively slightly wrong notation in correct arguments. \\

   \textbf{Misuse of notation:} wrong notation that leads to an incorrect argument.
\end{tcolorbox}

\begin{addmargin}[0.6cm]{0.6cm}  
\begin{center}
        \textit{\textcolor{airforceblue}{The next proposition is a univariate prototype of what we are going to develop in the case of $n$ variables. The proposition explains gcd in terms of ideals and it 
        addresses the question on how to do computation with ideals in the ring $k[x]$.}}
    \end{center}
\end{addmargin}

\begin{proposition}[title = Proposition 1.5.7]
    Let $f, g \in k[x]$. Then the following hold:
    \begin{itemize}
      \item[(i)] $gcd(f,g)$ exists and is unique up to a non-zero constant multiple (i.e. up to a factor in $k\backslash \{0\}$). 
      \item[(ii)] $gcd(f,g)$ generate the ideal $\langle f,g \rangle$, meaning that $\langle gcd(f,g) \rangle = \langle f,g \rangle$.   \\
            \hspace*{0.5cm}\textit{\textcolor{airforceblue}{(This is an equivalent descrtiption of the gcd in the language of ideals.)}}
      \item[(iii)] There exists an algorithm that determines $gcd(f, g)$ for given $f$ and $g$. 
    \end{itemize} 

    \paragraph{Proof:} Note that $f$ and $g$ may both be 0. WE have not excluded this possibility in our assumptions. This is a boundary case we deal with seperately. 

    Note that $f$ being divisible by $h$ means that $f = h \cdot g$ for some $g \in k[x]$. This also applies to the degenerate situations. For example, $0 = 0 \cdot 1$ means that 
    0 is divisible by 0. if we keep this in mind, we see that $gcd(0,0) = 0$ as 0 is divisible by 0 and every polynomial divides 0.
    
    ON the other hand, if $f$ or $g$ is not 0, then $gcd(f,g) \neq 0$. Say, if $f \neq 0$, then $f$ is not divisible by zero as $f \neq q  0$ for every $g \in k[x]$, which implies 
    that 0 is not a common divisor of $f$ and $g$. The above shows that \[gcd(f,g) = 0 \Leftrightarrow f = g =0\]  

    In this degenerate case all these three assertions (i), (ii) and (ii) are true.  \\

    So, in what follows we may assume $gcd(f,g) \neq 0$. The plan for the proof is to find a polynomial $h$ satisfying $\langle h \rangle = \langle f,g \rangle$ and then to 
    verify that this polynomial is the $gcd(f,g)$ obtaining (ii).ONce this is settled we will show (i) and eventually derive an algorithm for computing $gcd(f,g)$ to extablish (iii). \\

    Corollary 1.5.4 tells us that every ideal $i \subseteq k[x]$ can be generated by just one polynomial. Applying this to $I = \langle f,g \rangle$, we see that $\langle h \rangle = \langle f,g \rangle$ 
    for some $h \in k[x]$ without yet knowing how $h$ is related to $f$ and $g$ in a more elementary way (via gcd) and without knowing how to construct $h$ arithmetically. 
    Let's see why $h$ is $gcd(f,g)$. Before we see that $h$ is the \underline{greatest} common divisor of $f$ and $g$, we first check that it is a common divisor.  \\

    \begin{itemize}
      \item[(ii)] The set equality $\langle h \rangle -= \langle f,g \rangle$ is two inclusions: $\langle h \rangle \subseteq \langle f,g \rangle$ and $\langle h \rangle \supseteq \langle f,g \rangle$. \\[0.3em]
      The inclusion $\langle h \rangle \supseteq \langle f,g \rangle$ implies that $f \in \langle h \rangle$ and $g \in \langle h \rangle$. Hence, $h$ is a multiple of $f$ and 
      a multiple $g$. Thus, $h$ is indeed a common divisor of $f$ and $g$. 
      
      Now consider any common divisor $p$ of $f$ and $g$. Wehave $f = C \cdot p, \ g = D \cdot p$ for some $C,D \in k[x]$. \\[0.3em] 
      Inclusion $\langle h \rangle = \langle f, g \rangle$ tells us that $h \in \langle f,g \rangle$, meaning that $h = A \cdot f + B \cdot g$ for some $A,B \in k[x]$. \\[0.3em]  
      Putting things together we have \[h = A \cdot f = B \cdot g = A \cdot C \cdot p + B \cdot D \cdot p = (A \cdot C + B \cdot D) \cdot p\] 
      Thus, $h$ is divisible by $p$. If follows that $h$ is the \underline{greatest} common divisor of $f$ and $g$. 

      \item[(i)] Why uniqueness up to a non-zero constant multiple? 
      
      Remember that we assume $gcd(f,g) \neq 0$ as the case $gcd(f,g) = 0$ is already done. If $h$ and $H$ are both greatest common divisors, then $h$ divides $H$ and $H$ divides $h$. 
      As $gcd(f,g) \neq 0$, both are non-zero polynomials. So, $h$ divides $H$ implies $deg \ h \leq deg \ H$ and $H$ divides $h$ implies $deg \ H \leq deg \ h$. So, $deg \ h  = deg \ H$. 
      $H$ begin divisible by $h$ implies $H = h \cdot q$ for some $g \in k[x]$. But then $deg \ H = deg(h \cdot g) = deg \ h + deg \ g$. 

      In viel of $deg \ H = deg \ h$, we obtain $deg \ q =0$, meaning that $g \in k \backslash\{0\}$. This conclues the poof of (i). 
    \end{itemize}
  \end{proposition}

\clearpage

\begin{proposition}[title = Proposition 1.5.7 (continued)]
    \begin{itemize}
      \item[(iii)] We are going to describe the standard algorithm for computing $gcd(f,g)$, the so called \textit{Euclidean Algorithm}. It is an iterative algorithm that uses long division in each iteration. \\[0.3em]
      $gcd(f, 0) = f$, when $g = 0$, we are done. 
      If however $g \neq 0$, we can divide $f$ by $g$ obtaining the quotient $g \in k[x]$ and the remainder $r \in k[x]$ satisfying $f = q \cdot g + r, \ deg r < deg g$.
      Using $q, g, r$ out $gcd(f, g)$ is now expressed as $gcd(f ,g) = gcd(q \cdot g + r, g)$. We claim that $gcd(q \cdot g + r, g) = gcd(g, r)$. \\[0.3em] 
      Why is this true? 

      Every polynomial $p$ that divides $r$ and $g$ also divides $r = (q \cdot g + r) - q \cdot g$ and $g$. It follow that the pairs $(q \cdot g + r, g)$ and $(r, g)$ have the 
      same set of common divisors. Thus, they also have the same \underline{greatest} common divisor. We have $gcd(f,g) = gcd(g, r)$/ This observation can be used to set up an iterative 
      procedure: 

      $g$ is going to become the new $f$ and $r$ is the new $g$. During the iteration, the pair of the polynomials being processed will change, but the gcd of the pair will stay the same 
      while the degree of the second polynomial of the pair will rop with each iteration (because $deg \ g > deg \ r$) until the case $g = 0$ is reached. 

      Here is the complete algorithm: 

      \begin{algorithm}[title = Greatest Common Divisor]
        \textbf{Assumption}: $f,g \in k[x]$  \\
        \textbf{Result}: $gcd(f,g)$ 
\begin{lstpseudo}
% Invariant of the while loop: 
    gcd(f,g) does not change 

while g != 0: 
    q,r := Long_Division(f,g) 
    f := g 
    g := r 
    % deg g has dropped: 
        we are making progress, see you in the next iteration :) 
end while 

return f % gcd(f,g) = f when g = 0
\end{lstpseudo}
    \end{algorithm}

      There is not much to be sair on why this algorithm correctly determines $gcd(f,g)$ after finitely many iterations. The discussion before the algorithm and the comment lines 
      explain it all.
    \end{itemize}
      \begin{flushright}
        $\Box $
      \end{flushright}
\end{proposition} 

\begin{addmargin}[0.6cm]{0.6cm}  
\begin{center}
        \textit{\textcolor{airforceblue}{The case of an ideal $\langle f,g \rangle$ generated by two polynomials $f$ and $g$ in $k[x]$ is the base case for the case of an ideal 
        $\langle f_1, \ldots, f_s \rangle$ generated by finitely many polynomials $f_1, \ldots, f_s \in k[x]$. In this more general situation gcd is still the concept that \enquote{explains everything}.}}
    \end{center} 
\end{addmargin} 

\begin{definition}[title = Definition 1.5.8]
  Greatest common divisor of finitely many polynomials $f_1, \ldots, f_s \in k[x]$ is a polynomial $h \in k[x]$ such that: 
  \begin{itemize}
    \item[(i)] $h$ divides all of the $f_1, \ldots, f_s$ (\textit{\textcolor{airforceblue}{it is a \underline{common} divisor}}) 
    \item[(ii)] Every polynomial that divides all of the $f_1, \ldots, f_s$ also divides $h$ (\textit{\textcolor{airforceblue}{it is the \underline{greatest} of all common divisors}})  
  \end{itemize}
  Notation: $h = gcd(f_1, \ldots, f_s)$
\end{definition}

Next, we give a proposition that explains how to calculate  $gcd(f_1, \ldots, f_S)$ for $s \geq 2$ iteratively using the base case $s = 2$, which we have already done 
and which gives the iterpretation of $gcd(f_1, \ldots f_s)$ in term of ideals. The proposition aso clarifies uniqueness. So, we have all of the properties of the claims of 
Proposition 1.5.7 extendet to the case of $s \geq 2$ polynomials. 

\begin{proposition}[title = Proposition 1.5.9]
    For $s \geq 2$ polynomials $f_1, \ldots, f_s \in k[x]$, the following assertions hold:
    \begin{itemize}
      \item[(i)] $gcd(f_1, \ldots, f_s)$ is a unique cap to a non-zero multiple from $k\backslash\{0\}$. 
      \item[(ii)] $\langle gcd(f_1, \ldots, f_s) \rangle = \langle f_1, \ldots, f_s \rangle$\\
      \textit{\textcolor{airforceblue}{(Thas is $gcd(f_1, \ldots, f_s)$ is characterized as the generator of the ideal $\langle f_1, \ldots, f_s \rangle$)}}
      \item[(iii)] $gcd(f_1, \ldots, f_s) = gcd(f_1, gcd(f_2, \ldots, f_s))$ 
      \item[(iv)] There exists an algorithm that calculates $gcd(f_1, \ldots, f_s)$ for given $f_1, \ldots, f_s \in k[x]$.  
    \end{itemize} 

    \paragraph{Proof:} To show (i) and (ii) just repeat the arguments in the proof of assertions (i) an (ii) of Proposition 1.5.7. Those assertions were for two polynomials. 
    Now, we have $s$ polynomials buXXt the agruments are completely the same. 
    \begin{itemize}
      \item[(iii)] can be verified either directly or as a consequence of (ii). Verifying using (ii) is probably the more interesting way to do it. 
      
      We know $\langle \gcd(f_2, \ldots, f_s) \rangle = \langle f_2, \ldots, f_s \rangle$ (*). So, 
      \begin{align*}
        \langle f_1, ..., f_s \rangle 
        & = \{f_1 \cdot h_1 + \underbrace{f_2 \cdot h_2 \ldots + f_s \cdot h_s}_{
          \mathclap{
            \substack{
                \text{ can be replaced by a multiple of }
              \\
                \gcd(f_2, \ldots, f_S) \text{ in a view of (*)}
          }}} \mid h_1, \ldots, h_s \in k[x]
        \} \\
        & = \langle f_1, gcd(f_2, \ldots, f_s) \rangle \stackrel{(*)}{=}  \langle gcd(f_1, gcd(f_2, \ldots, f_s)) \rangle \\ 
        \intertext{ (*) we use the description of the gcd in terms of ideals given in (ii)}
        \langle f_1, \ldots, f_s \rangle &= \langle \underbrace{gcd(f_1, gcd(f_2, \ldots, f_s))}_{
          \mathclap{
            \substack{
              \text{this is} \gcd(f_1, \ldots,f_s) \\ \text{ by the description of gcd in terms} \\ text{of ideals, given in (ii).}
            }
          }
        } \rangle
      \end{align*}
    \end{itemize}

    In view of (iii), we can verify (iv) by showing that calculation of the gcd of $s$ polynomials reduces to the iterative process, in which the gcd of \underline{two} polynomials 
    is calculated in each iteration.  

    \begin{algorithm}[title = Greatest Common Divisor]
        \textbf{Assumption}: $f,g \in k[x]$  \\
        \textbf{Result}: $gcd(f,g)$ 
\begin{lstpseudo}
i := s 
g :=f_s   % this is the gcd in progress 

% invariant of the loop: g = gcd(f_i, ..., f_s)

while i > l:  % work to be done? 
    g := greatest_common_divisor(f_{i-1}, g) 
    % not f_{i-1} is also accounted for 
    i := i - 1  % uodate i to maintain the invariant in the next iteration 
end while 

return g   % as i = 1, this is the answer

\end{lstpseudo}
    \end{algorithm}

\begin{addmargin}[0.6cm]{0.6cm}  
    \begin{center}
            \textit{\textcolor{airforceblue}{There is an alternative approach to solve this algorithmic problem. We could have extended the Euclidean Algorithm from two polynomials to $s \geq 2$ polynomials.}}
    \end{center} 
\end{addmargin}

      \begin{flushright}
        $\Box $
      \end{flushright}
\end{proposition} 

\begin{addmargin}[0.6cm]{0.6cm}  
    \begin{center}
            \textit{\textcolor{airforceblue}{With these result we are ready to solve the ideal membership problem for the ring $k[x]$:}} \\
            \textit{\textcolor{airforceblue}{deciding if a given polynomial belongs to a given ideal $\langle f_1, \ldots, f_s \rangle$ when $f_1, \ldots, f_S \in k[x]$.}}
    \end{center} 
\end{addmargin}

\clearpage
\begin{corollary}[title = Corollary 1.5.10]
    There is an algorithm that decides of a polynomial $f$ is in the ideal $\langle f_1, \ldots, f_s \rangle$ for given $f$ and $f_1, \ldots, f_s \in k[x]$.

    \paragraph{Proof:} Calculate the $\gcd(f_1, \ldots, f_s)$ and check if $f$ is divisible by $gcd(f_1, \ldots, f_s)$. Since $\langle f_1, \ldots, f_s \rangle = \langle \gcd(f_1, \ldots, f_s) \rangle$ 
    (by Proposition 1.5.9 (ii)) and since $\gcd(f_1, \ldots, f_s)$ can be computed (by Proposition 1.5.9 (iv)), this procedure solves the decision problem \[f \in \langle f_1, , f_s \rangle.\]

      \begin{flushright}
        $\Box $
      \end{flushright}
\end{corollary}

\begin{exercise}[title = Exercise 1.5.11]
  \begin{addmargin}[0.6cm]{0.6cm}  
    \begin{center}
            \textit{\textcolor{airforceblue}{This section had algorithms so that we can have computational exercises just to tryy the algorithms out. Otherwise 
            we can see how these algorithms help solving systems of polynomial equations in one variable when the underlying coefficnet field is algebraically closed. 
            Concretely, we consider the field $k = \mathbb{C}$} of complex numbers.} \\
    \end{center} 
  \end{addmargin} 

  \begin{enumerate}
    \item Carry out long division of $x^3 + 2x_2 + x + 1$ by $2x+ 1$ in $\mathbb{Q}[x]$. 
    \item Determine $\gcd(x^4 -1, x^6 -1)$ in $\mathbb{Q}[x]$. 
    \item Show that every polynomial $ f \in \mathbb{C}[x]$ of degree $n \geq 1$ can be factoriced as $f = c (x-r_1)^{m_1} \cdot \ldots \cdot (x - r_l)^{m_l}$ where $c \in \mathbb{C}\backslash \{0\}, r_1, \ldots, r_l$ are distinct coots of $f$, 
    $m_1, \ldots, m_l \geq 1$ and $m_1 + \ldots + m_l = n$. USe the fact, that $\mathbb{C}$ is algebraically closed. 
    \item \textit{\textcolor{airforceblue}{This one is a follow-up on the Euclidean Algorithm.}} \\
    Show that for $f,g \in k[x]$ ine has \[\underbrace{gcd(f,g) = A \cdot f+ B  \cdot g}_{\text{Bézout's identity}}\] for some $A,B \in k[x]$.  \\
    \textit{\textcolor{airforceblue}{KEywords: Extended Euclidean Algorithm (for polynomials)}}
    \item Is $x^3 + 4x^2 + 3x + 7 \in \langle x^3-3x + 2, x^4 -1, x^6 -1 \rangle$? The ring here is $\mathbb{Q}[x]$. 
    \item Consider the system of polynomial equations $f_1 = \ldots = f_s = 0$ in one variable $x$ an with coefficients in $\mathbb{C}$. Design an algorithm that verifies 
    if this system has a solution in $\mathbb{C}$. The arithmetic of $\mathbb{C}$ is assumed to be available. 
    \item \begin{addmargin}[0.6cm]{0.6cm}  
    \begin{center}
            \textit{\textcolor{airforceblue}{You must have had some experience with calculus in you Bachelor program. So, i guess you know what a derivative is.
            While one cannot do derivatives over any field the same way as one does them for the fields $k = \mathbb{R}$ or $k = \mathbb{C}$ in calculus, one can 
            define the derivative $f'$ if a polynomial $f \in k[x]$ for \underline{any} coefficient field $k$ completely formally:}} \\
 
            \textit{\textcolor{airforceblue}{If $f = \sum_{j = 0}^{\infty} c_j x^j $}, then $f'$ is defined by $f' = \sum_{j = 1}^{\infty} j \cdot c_j x^{j-1}.$}
    \end{center} 
  \end{addmargin} 
    Show 
    \begin{align*}
      (f + g)' & f' + g' \\ 
      (f \cdot g)' & f' \cdot g + f \cdot g' 
    \end{align*}

    \item Assume that a polynomial $f \in \mathbb{C}[x]$ of degree $\geq 1$ has been factorized as in 3: \[f = c(x-r_1)^{m_1} \cdot \ldots \cdot (x - r_l)^{m_l}\]
    Find a representive factorization 0f $\gcd(f, f')$ and $\frac{g}{\gcd(f, f')}$.  

    \item 
    \begin{center}
            \textit{\textcolor{airforceblue}{If you have solved the above exercises, you probably understand that the system $f_1= \ldots = f_s = 0$ with 
            $f_1, \ldots, f_s \in k[x]$ gets reduced to just one equation $f = 0$. Its solution set is the variety $V(f)$.}}  \\
     \end{center} 
     Show for $k = \mathbb{C}$, that \[I(v(f)) = \langle \frac{f}{gcd(f, f')} \rangle\]
  \end{enumerate} 
\end{exercise}